% from ega4-11-fr.tex
\oldpage[IV-3]{141}

comme $X'_\alpha$ et $Y'_\alpha$ sont noethériens et que $A''_\alpha$ est
intégralement clos, on déduit de (6.5.4) que $X'_\alpha$ est normal au point
$x'_\alpha$. Mais les morphismes
$\operatorname{Spec}(A''_\beta)\to\operatorname{Spec}(A''_\alpha)$ pour
$\alpha\leq\beta$ sont dominants, donc, comme en vertu de (11.3.1) on peut
supposer que les $f'_\alpha$ sont plats pour $\alpha\geq\alpha_0$, il résulte
de (2.3.7) que toute composante irréductible de $X'_\beta$ domine une
composante irréductible de $X'_\alpha$. On conclut alors de (5.13.4), appliqué
aux préschémas
\[
\operatorname{Spec}(\mathcal O_{X_0,x_0}\otimes_{A_0}A''_\alpha),
\]
que $X$ est normal au point $x$.
\end{proof}

\begin{corollary}[11.3.14]
\label{IV.11.3.14-fr}
Soient $f:X\to Y$ un morphisme localement de présentation finie, $x$ un point
de $X$, $y=f(x)$. Si $Y$ est géométriquement unibranche au point $y$ (6.15.1)
et si le morphisme $f$ est normal au point $x$ (6.8.1), alors $X$ est
géométriquement unibranche au point $x$.
\end{corollary}

\begin{proof}
On peut évidemment se borner au cas où $Y=\operatorname{Spec}(A)$, $A$ étant
un anneau local intègre géométriquement unibranche, $y$ étant le point fermé de
$Y$. Soit $A'$ la clôture intégrale de $A$~; posons
$Y'=\operatorname{Spec}(A')$ et soit $y'$ le point fermé de $Y'$, de sorte que
le morphisme $g:Y'\to Y$ est radiciel au point $y$ (6.15.3), entier et
birationnel. Si l'on pose $X'=X\times_Y Y'$, et si $f':X'\to Y'$ et
$g':X'\to X$ sont les projections, $g'$ est entier et radiciel au point $x$
(6.15.3.1). Soit donc $x'$ l'unique point de $g'^{-1}(x)$, qui est au-dessus
de $y'$. Le morphisme $f'$ est de présentation finie et normal au point $x'$
(6.7.8), et par suite l'anneau local $\mathcal O_{X',x'}$ est intègre et
intégralement clos (11.3.13). D'autre part, on peut supposer $f$ plat (11.3.1),
donc $g'$ est un morphisme birationnel (6.15.4.1)~; par suite
$\operatorname{Spec}(\mathcal O_{X,x})$ est irréductible, et comme
$\mathcal O_{X,x}$ est réduit (11.3.13) il est intègre et géométriquement
unibranche en vertu de (6.15.5).
\end{proof}

\begin{proposition}[11.3.15]
\label{IV.11.3.15-fr}
Soient $A$ un anneau, $B$ une $A$-algèbre de présentation finie, $M$ un
$B$-module de présentation finie, qui est un $A$-module plat. Alors il existe
une suite finie $(f_i)_{1\leq i\leq n}$ d'éléments de $A$, tels que l'idéal
engendré par les $f_i$ soit égal à $A$, et que, pour $0\leq i<n$,
\[
M_{f_{i+1}}\Big/\Big(\sum_{j=1}^{i}f_jM_{f_{i+1}}\Big)
\]
soit un
$A_{f_{i+1}}\big/(\sum_{j=1}^{i}f_jA_{f_{i+1}})$-module libre.
\end{proposition}

On peut encore dire que les $D(f_i)$ forment un recouvrement ouvert de
$X=\operatorname{Spec}(A)$, et que si l'on pose $Z_1=D(f_1)$, puis, par
récurrence,
\[
Z_{i+1}=D(f_{i+1})\cap V(f_1A+\cdots+f_iA),
\]
les $Z_i$ forment une \emph{partition} de $X$ en ensembles localement fermés,
$A_{f_{i+1}}/(\sum_{j=1}^{i}f_jA_{f_{i+1}})$ étant l'anneau d'un sous-schéma
affine de $X$ ayant $Z_{i+1}$ pour espace sous-jacent.

\begin{proof}
Pour démontrer la proposition, on peut d'abord, en vertu de (11.2.7), supposer
qu'il existe un sous-anneau noethérien $A_0$ de $A$, une $A_0$-algèbre de type
fini $B_0$ et un $B_0$-module de type fini $M_0$, plat sur $A_0$ et tels que
$B=B_0\otimes_{A_0}A$, $M=M_0\otimes_{A_0}A$~; il est clair qu'il suffira de
prouver la proposition pour $A_0$, $B_0$ et $M_0$, car si les éléments
$f_i\in A_0$ vérifient dans ce cas les conditions de l'énoncé, ils les
vérifieront aussi pour $A$, $B$, $M$, puisque
\[
M_{f_{i+1}}\Big/\Big(\sum_{j=1}^{i}f_jM_{f_{i+1}}\Big)
=\left((M_0)_{f_{i+1}}\Big/\Big(\sum_{j=1}^{i}f_j(M_0)_{f_{i+1}}\Big)\right)
\otimes_{A_0}A
\]
(Bourbaki, \emph{Alg. comm.}, chap.~II, \S~2, n$^\circ$~7, prop.~18). On
peut donc se borner désormais au cas où $A$ est noethérien.

Notons maintenant que si $C$ est un anneau noethérien, $\mathfrak N$ son
nilradical et $P$ un $C$-module plat, alors il résulte de
$(\mathbf 0_{\mathrm{III}},10.1.2)$ que pour que $P$ soit un $C$-module libre,

